\Chapter{Architektura i technologie}
\label{ch:architecture}

\section{Podział odpowiedzialności programu}

Projekt jest pakietem Pythona w układzie \code{src/pitstop}.
Wspólnym punktem wejścia jest polecenie \code{pitstop}, zbudowane za
pomocą \code{argparse}. Oddzielne moduły odpowiadają etapom przetwarzania,
a \code{contracts.py} definiuje nazwy, typy, kolejność cech i walidatory.
Taki podział ogranicza zależność między autorami kolejnych etapów:
odbiorca może sprawdzić format wejścia bez znajomości sposobu jego
wytworzenia. Sam poprawny format nie dowodzi jednak poprawności
merytorycznej etykiety ani przyczynowości cech.

\begin{figure}[htbp]
\centering\small
\begin{tabular}{c}
\fbox{\strut Człowiek: pobranie do cache FastF1} \\[5pt]
$\downarrow$ \\[3pt]
\fbox{\strut T2 offline: normalizacja $\to$ snapshot i manifest $\to$ weryfikacja} \\[5pt]
$\downarrow$ \\[3pt]
\fbox{\strut Plan: etykiety $\to$ podziały $\to$ cechy i okna $\to$ zbiór} \\[5pt]
$\downarrow$ \\[3pt]
\fbox{\strut Plan: uczenie $\to$ predykcje $\to$ ocena $\to$ panel}
\end{tabular}
\caption{Przepływ danych z zaznaczeniem działającej i planowanej części.
Opracowanie na podstawie modułów projektu.}
\label{fig:flow}
\end{figure}

CLI obsługuje eksport \code{download --from-cache} z wymaganymi
argumentami oraz weryfikację \code{verify-data --manifest}.
Nazwa pierwszego polecenia wynika ze wspólnego interfejsu, ale ten
wariant eksportuje istniejący cache bez dostępu do sieci.
Polecenia \code{build-dataset}, \code{train}, \code{predict},
\code{evaluate} i \code{export-demo} są zapowiedzią interfejsów:
obecnie wypisują komunikat o braku implementacji i kończą się kodem 1.
Moduły modeli zawierają opisy zamierzonej odpowiedzialności, bez klas
realizujących obliczenia~\cite{project}.

\section{Środowisko i rola bibliotek}

Projekt wymaga Pythona 3.12. Menedżer \code{uv} odtwarza środowisko
z \code{pyproject.toml} i wersjonowanego \code{uv.lock}.
Tabela~\ref{tab:tech} podaje wersje zapisane w tej blokadzie; nie jest
wykazem bibliotek użytych w przeprowadzonym treningu.

\begin{table}[htbp]
\centering\small
\caption{Technologie zapisane w źródłowej rewizji projektu.}
\label{tab:tech}
\begin{tabularx}{\textwidth}{L{0.24\textwidth}L{0.18\textwidth}Y}
\toprule
Technologia & Wersja & Rola i aktualny zakres \\
\midrule
FastF1 & 3.9.0a0 & Odczyt cache oraz audyt źródła. \\
pandas / NumPy & 2.3.3 / 2.5.4 & Tabele, typy i walidacja tablic. \\
PyArrow & 26.0.0 & Zapis i odczyt Parquet. \\
PyYAML & 6.0.3 & Odczyt kompletnych konfiguracji YAML. \\
scikit-learn & 1.9.1 & Zależność dla planowanych metod i metryk. \\
PyTorch & 2.14.1 & Opcjonalne \code{train}; planowane sieci i pętla uczenia. \\
LightGBM & 4.7.0 & Opcjonalne \code{train}; planowany model porównawczy. \\
pytest / pytest-socket & 9.1.1 / 0.8.1 & Testy i blokada gniazd sieciowych. \\
Ruff & 0.17.1 & Kontrola i formatowanie kodu. \\
\bottomrule
\end{tabularx}
\end{table}

Blokada zawiera także wariant \code{2.14.1+cpu} PyTorch dla właściwych
platform; CPU jest urządzeniem referencyjnym. CUDA i MPS są opcjonalne
w projekcie docelowym. Streamlit jest planowaną technologią panelu,
lecz w tej rewizji nie ma jeszcze aplikacji ani jej osobnego środowiska.
Nie można zainstalowanej biblioteki utożsamiać z gotowym komponentem.

Konfiguracje nie dziedziczą po sobie: każdy plik modelu powtarza potrzebne
ścieżki i parametry. Ułatwia to zapis pełnego opisu przyszłego eksperymentu
bez ukrytych ustawień. Ziarno losowości zaczyna się od 42, ale obecność
tej wartości w konfiguracji nie jest dowodem wykonania eksperymentu.

\section{Automatyczna kontrola zmian}

GitHub Actions uruchamia instalację z blokady, \code{ruff check},
\code{ruff format --check} i \code{pytest}. CI ustala \code{uv 0.11.7}
i Pythona 3.12. Pobranie zależności jest częścią przygotowania środowiska;
same testy mają blokadę połączeń internetowych. Nie korzystają z konta
F1 TV ani z pełnego cache wyścigowego.

D10 przewiduje w przyszłości także krótki trening na danych testowych,
ale aktualny plik workflow nie zawiera takiego kroku. Sieci, pętla uczenia
i panel wymagają dopiero implementacji. To rozróżnienie pozwala odczytać
zielony wynik CI zgodnie z rzeczywistym zakresem kontroli.
